\documentclass[10pt,a4paper]{amsart}
\usepackage{amsmath,amssymb,amsthm,graphicx}
% For commutative diagrams we use Xy-pic
\usepackage[all]{xy}

% We use 2cell for 2-commutative diagrams.
\xyoption{2cell}
% Theorem environments.
%
\theoremstyle{plain}
\newtheorem{theorem}[subsection]{Theorem}
\newtheorem{proposition}[subsection]{Proposition}
\newtheorem{lemma}[subsection]{Lemma}

\theoremstyle{definition}
\newtheorem{definition}[subsection]{Definition}
\newtheorem{example}[subsection]{Example}
\newtheorem{exercise}[subsection]{Exercise}
\newtheorem{situation}[subsection]{Situation}

\theoremstyle{remark}
\newtheorem{remark}[subsection]{Remark}
\newtheorem{remarks}[subsection]{Remarks}

\numberwithin{equation}{subsection}

% Macros
%
\def\lim{\mathop{\mathrm{lim}}\nolimits}
\def\colim{\mathop{\mathrm{colim}}\nolimits}
\def\Spec{\mathop{\mathrm{Spec}}}
\def\Hom{\mathop{\mathrm{Hom}}\nolimits}
\def\Ext{\mathop{\mathrm{Ext}}\nolimits}
\def\SheafHom{\mathop{\mathcal{H}\!\mathit{om}}\nolimits}
\def\SheafExt{\mathop{\mathcal{E}\!\mathit{xt}}\nolimits}
\def\Sch{\mathit{Sch}}
\def\Mor{\mathop{\mathrm{Mor}}\nolimits}
\def\Ob{\mathop{\mathrm{Ob}}\nolimits}
\def\Sh{\mathop{\mathit{Sh}}\nolimits}
\def\NL{\mathop{N\!L}\nolimits}
\def\CH{\mathop{\mathrm{CH}}\nolimits}
\def\proetale{{pro\text{-}\acute{e}tale}}
\def\etale{{\acute{e}tale}}
\def\QCoh{\mathit{QCoh}}
\def\Ker{\mathop{\mathrm{Ker}}}
\def\Im{\mathop{\mathrm{Im}}}
\def\Coker{\mathop{\mathrm{Coker}}}
\def\Coim{\mathop{\mathrm{Coim}}}

% Boxtimes
%
\DeclareMathSymbol{\boxtimes}{\mathbin}{AMSa}{"02}

%
% Macros for moduli stacks/spaces
%
\def\QCohstack{\mathcal{QC}\!\mathit{oh}}
\def\Cohstack{\mathcal{C}\!\mathit{oh}}
\def\Spacesstack{\mathcal{S}\!\mathit{paces}}
\def\Quotfunctor{\mathrm{Quot}}
\def\Hilbfunctor{\mathrm{Hilb}}
\def\Curvesstack{\mathcal{C}\!\mathit{urves}}
\def\Polarizedstack{\mathcal{P}\!\mathit{olarized}}
\def\Complexesstack{\mathcal{C}\!\mathit{omplexes}}
% \Pic is the operator that assigns to X its picard group, usage \Pic(X)
% \Picardstack_{X/B} denotes the Picard stack of X over B
% \Picardfunctor_{X/B} denotes the Picard functor of X over B
\def\Pic{\mathop{\mathrm{Pic}}\nolimits}
\def\Picardstack{\mathcal{P}\!\mathit{ic}}
\def\Picardfunctor{\mathrm{Pic}}
\def\Deformationcategory{\mathcal{D}\!\mathit{ef}}
\usepackage{cleveref}
\crefname{originaldefinition}{original definition}{original definitions}
\Crefname{originaldefinition}{Original definition}{Original definitions}
\makeatletter
\newlabel{situation-setup}{{02QL}{source}{}{sourcepointer-1}{}}
\newlabel{situation-setup@cref}{{[originaldefinition][0][]02QL}{[1][0][]source}{}{}{}}
\makeatother
\crefname{originallemma}{original lemma}{original lemmas}
\Crefname{originallemma}{Original lemma}{Original lemmas}
\makeatletter
\newlabel{lemma-from-chow-to-K}{{0FDS}{source}{}{sourcepointer-2}{}}
\newlabel{lemma-from-chow-to-K@cref}{{[originallemma][0][]0FDS}{[1][0][]source}{}{}{}}
\makeatother
\crefname{originallemma}{original lemma}{original lemmas}
\Crefname{originallemma}{Original lemma}{Original lemmas}
\makeatletter
\newlabel{lemma-cycles-k-group}{{02S9}{source}{}{sourcepointer-3}{}}
\newlabel{lemma-cycles-k-group@cref}{{[originallemma][0][]02S9}{[1][0][]source}{}{}{}}
\makeatother
\crefname{originallemma}{original lemma}{original lemmas}
\Crefname{originallemma}{Original lemma}{Original lemmas}
\makeatletter
\newlabel{homology-lemma-serre-subcategory-K-groups}{{02MX}{source}{}{sourcepointer-4}{}}
\newlabel{homology-lemma-serre-subcategory-K-groups@cref}{{[originallemma][0][]02MX}{[1][0][]source}{}{}{}}
\makeatother
\crefname{originallemma}{original lemma}{original lemmas}
\Crefname{originallemma}{Original lemma}{Original lemmas}
\makeatletter
\newlabel{lemma-maps-between-coherent-sheaves}{{02SB}{source}{}{sourcepointer-5}{}}
\newlabel{lemma-maps-between-coherent-sheaves@cref}{{[originallemma][0][]02SB}{[1][0][]source}{}{}{}}
\makeatother
\title{The Chow group in dimension k is the image of the K-group of sheaves supported in dimension at most k inside the K-group of sheaves supported in dimensions k+1 and k}
\author{The Stacks Project Authors}
\date{stacks-project:chow.tex:local-snapshot-20261012; GNU FDL 1.2 or later}
\begin{document}
\maketitle
\noindent\textit{Editorial note.} All delivered mathematical text is the author's own, copied from the pinned Stacks Project chapter file chow.tex; no statement or proof step was rewritten, reordered or removed.\par\medskip
\noindent\textit{Editorial note.} Preamble provenance: the pinned chapter file opens with an input of the shared preamble and carries no preamble text of its own, so the project's single Stacks-wide preamble staged at sources/stacks-chow/preamble.tex is excerpted verbatim as two spans, lines 14-18 (Xy-pic setup) and lines 154-217 (the native theorem declarations and mathematical macros). Nothing was added, removed or redefined; the wrapper class and the amsmath/amssymb/amsthm/graphicx base packages are supplied by the assembler.\par\medskip
\noindent\textit{Editorial note.} The counted statement and proof contain no cross reference, no citation and no included file, so the delivered body is self contained.\par\medskip
\noindent\textit{Editorial note.} Canonical location of the counted proposition: Stacks tag 02SC (sources/stacks-chow/pinned-original-source.tex line 18509).\par\medskip
\noindent\textit{Allowed standard prerequisites (editorial).}\begin{itemize}
\item Identification of Z\_k(X) with the K-group of the quotient of the category of sheaves supported in dimension at most k. Same chapter, Lemma labelled lemma-cycles-k-group (source line 3757) and the construction of Lemma labelled lemma-from-chow-to-K (source line 3934), invoked at statement lines 18517 and proof lines 18527-18529. Fix the map whose bijectivity modulo the subgroup B\_k(X) is the content of the lemma, and identify the source with the cycle group.
\item K-groups of a Serre subcategory and its quotient: killing a class and realising it by a (2, 1)-periodic complex. Homology, Lemma labelled homology-lemma-serre-subcategory-K-groups (Stacks tag 02MX), parts (1) and (2), invoked at proof lines 18545-18549. Supplies both the passage from the vanishing of [A] - [B] in the larger quotient to the existence of A', B' supported in dimension at most k-1 and the realisation of the resulting class by a (2, 1)-periodic complex; the cycle comparison is done in the counted span.
\item Rational equivalence of the k-cycles of the two cohomology sheaves of such a complex. Same chapter, Lemma labelled lemma-maps-between-coherent-sheaves (source line 18417), invoked at proof line 18554. Supplies [A + A']\_k \textasciitilde{} rat [B + B']\_k for a (2, 1)-periodic complex satisfying the support conditions; the reduction of the K-theoretic statement to this input is inside the counted span.
\item The base situation (S, delta) and the dimension conventions for supports. Same chapter, Situation labelled situation-setup (source line 1372), invoked in the statement. Fixes the delta-dimension, the categories Coh\_\{<= k\}(X) and the cycle group Z\_k(X) in which the statement is formulated.
\end{itemize}\medskip
\ifdefined\hypertarget\hypertarget{sourcepointer-1}{}\fi
\noindent\textit{Source reference.} Original definition 02QL: sources/stacks-chow/pinned-original-source.tex, source lines 1372--1372. Same-chapter pointer into the author environment situation at source line 1372; invoked at proof line None of the counted proof. It is retained only as a navigation pointer; every load-bearing step is inside the delivered span.\par\medskip
\ifdefined\hypertarget\hypertarget{sourcepointer-2}{}\fi
\noindent\textit{Source reference.} Original lemma 0FDS: sources/stacks-chow/pinned-original-source.tex, source lines 3934--3934. Same-chapter pointer into the author environment lemma at source line 3934; invoked at proof line 18529 of the counted proof. It is retained only as a navigation pointer; every load-bearing step is inside the delivered span.\par\medskip
\ifdefined\hypertarget\hypertarget{sourcepointer-3}{}\fi
\noindent\textit{Source reference.} Original lemma 02S9: sources/stacks-chow/pinned-original-source.tex, source lines 3757--3757. Same-chapter pointer into the author environment lemma at source line 3757; invoked at proof line 18527 of the counted proof. It is retained only as a navigation pointer; every load-bearing step is inside the delivered span.\par\medskip
\ifdefined\hypertarget\hypertarget{sourcepointer-4}{}\fi
\noindent\textit{Source reference.} Original lemma 02MX: sources/stacks-chow/tags.txt, source lines 3460--3460. Cross-chapter standard prerequisite located in the project tag list staged in the local snapshot; invoked at proof line 18546 of the counted proof. It is retained only as a navigation pointer; every load-bearing step is inside the delivered span.\par\medskip
\ifdefined\hypertarget\hypertarget{sourcepointer-5}{}\fi
\noindent\textit{Source reference.} Original lemma 02SB: sources/stacks-chow/pinned-original-source.tex, source lines 18417--18417. Same-chapter pointer into the author environment lemma at source line 18417; invoked at proof line 18554 of the counted proof. It is retained only as a navigation pointer; every load-bearing step is inside the delivered span.\par\medskip
\begin{lemma}
\label{lemma-cycles-rational-equivalence-K-group}
Let $(S, \delta)$ be as in Situation \ref{situation-setup}.
Let $X$ be a scheme locally of finite type over $S$.
The map
$$
\CH_k(X) \longrightarrow
K_0(\textit{Coh}_{\leq k + 1}(X)/\textit{Coh}_{\leq k - 1}(X))
$$
from Lemma \ref{lemma-from-chow-to-K} induces a bijection from
$\CH_k(X)$ onto the image $B_k(X)$ of the map
$$
K_0(\textit{Coh}_{\leq k}(X)/\textit{Coh}_{\leq k - 1}(X))
\longrightarrow
K_0(\textit{Coh}_{\leq k + 1}(X)/\textit{Coh}_{\leq k - 1}(X)).
$$
\end{lemma}

\begin{proof}
By Lemma \ref{lemma-cycles-k-group} we have
$Z_k(X) = K_0(\textit{Coh}_{\leq k}(X)/\textit{Coh}_{\leq k - 1}(X))$
compatible with the map of Lemma \ref{lemma-from-chow-to-K}.
Thus, suppose we have an element $[A] - [B]$ of
$K_0(\textit{Coh}_{\leq k}(X)/\textit{Coh}_{\leq k - 1}(X))$
which maps to zero in $B_k(X)$, i.e., maps to zero in
$K_0(\textit{Coh}_{\leq k + 1}(X)/\textit{Coh}_{\leq k - 1}(X))$.
We have to show that $[A] - [B]$ corresponds to a cycle
rationally equivalent to zero on $X$.
Suppose $[A] = [\mathcal{A}]$ and $[B] = [\mathcal{B}]$
for some coherent sheaves $\mathcal{A}, \mathcal{B}$ on
$X$ supported in $\delta$-dimension $\leq k$.
The assumption that $[A] - [B]$ maps to zero in the group
$K_0(\textit{Coh}_{\leq k + 1}(X)/\textit{Coh}_{\leq k - 1}(X))$
means that there exists coherent sheaves
$\mathcal{A}', \mathcal{B}'$ on $X$ supported in
$\delta$-dimension $\leq k - 1$ such that
$[\mathcal{A} \oplus \mathcal{A}'] - [\mathcal{B} \oplus \mathcal{B}']$
is zero in $K_0(\textit{Coh}_{k + 1}(X))$ (use part (1) of
Homology, Lemma \ref{homology-lemma-serre-subcategory-K-groups}).
By part (2) of
Homology, Lemma \ref{homology-lemma-serre-subcategory-K-groups}
this means there exists a $(2, 1)$-periodic complex
$(\mathcal{F}, \varphi, \psi)$ in the category $\textit{Coh}_{\leq k + 1}(X)$
such that
$\mathcal{A} \oplus \mathcal{A}' = H^0(\mathcal{F}, \varphi, \psi)$
and $\mathcal{B} \oplus \mathcal{B}' = H^1(\mathcal{F}, \varphi, \psi)$.
By Lemma \ref{lemma-maps-between-coherent-sheaves}
this implies that
$$
[\mathcal{A} \oplus \mathcal{A}']_k
\sim_{rat}
[\mathcal{B} \oplus \mathcal{B}']_k
$$
This proves that $[A] - [B]$ maps to a cycle rationally
equivalent to zero by the map
$$
K_0(\textit{Coh}_{\leq k}(X)/\textit{Coh}_{\leq k - 1}(X))
\longrightarrow Z_k(X)
$$
of Lemma \ref{lemma-cycles-k-group}. This is what we
had to prove and the proof is complete.
\end{proof}

\end{document}
